\documentclass[exercice]{thedocument}%
\startdatakeys
\source{23GENMATATG11}
\cycle{4}
\competencesDuSocle{CA,RA}
\connaissancesRequises{A2a*,A2d*}
\competencesTravaillees{A2b*,A2e1,A2g*}
\enddatakeys
\begin{document}%
    Pour constituer les lots, on dispose de 195 figurines et 234 autocollants.
    Chaque lot sera composé de figurines ainsi que d'autocollants. Tous les
    lots sont identiques. Toutes les figurines et tous les autocollants doivent
    être utilisés.
    \begin{enumerate}
        \item Peut-on faire 3 lots ?  
        \item Décomposer 195 en produit de facteurs premiers. 
        \item Sachant que la décomposition en produit de facteurs premiers de 234
        est $2\times3^2\times13$~: 
        \begin{enumerate}
            \item Combien de lots peut-on constituer au maximum~? 
            \item De combien de figurines et d'autocollants sera alors composé
            chaque lot~?
        \end{enumerate}
    \end{enumerate}
\end{document}